\section{The trace-norm bound and every temperature in Hastings' thermal expansion}
\label{sec:hastings_thermal_trace_norm}

% Hastings, arXiv:cond-mat/0508554v3, proj3new.tex, lines 1281--1301 (the
% trace-norm form) and 1303--1320 (eq. mplt, every temperature).

Two steps of \cite[Section~IV]{Hastings2006Solving} follow the high-temperature
cluster expansion. The first turns its bound on expectation values into a bound
in trace norm. The second extends the matrix product form of
Section~\ref{sec:hastings_product_form} to every temperature by taking an even
power of the high-temperature form. This section proves both steps for general
operators. The truncated density matrix $\rho(\beta,l_{\mathrm{proj}})$ of the
cluster expansion enters only through the hypotheses, and the error bound at
every temperature \cite[lines 1322--1349]{Hastings2006Solving} is not proved
here. The trace norm of a square matrix $D$ is
$|D|_1=\operatorname{Tr}\sqrt{D^\dagger D}$, and $\|O\|$ is the operator norm.

\begin{theorem}[The trace norm from expectation values]
    \label{thm:hastings_thermal_trace_norm_duality}
    \lean{TNLean.Lattice.ClusterExpansion.re_trace_sqrt_le_of_forall_norm_trace_mul_le,
        TNLean.Lattice.ClusterExpansion.l2OpNorm_le_one_of_mem_unitaryGroup}
    \leanok
    Let $D$ be a square complex matrix and $\epsilon$ a real number. If
    $|\operatorname{Tr}(OD)|\le\epsilon$ for every matrix $O$ with
    $\|O\|\le1$, then $|D|_1\le\epsilon$.
\end{theorem}
\begin{proof}\leanok
    The polar decomposition writes $D=U\sqrt{D^\dagger D}$ with $U$ unitary.
    Then $O=U^\dagger$ has $\|O\|\le1$ and
    $\operatorname{Tr}(OD)=\operatorname{Tr}\sqrt{D^\dagger D}=|D|_1$. For a
    Hermitian $D$ this $O$ may be taken to be the operator of signs of the
    eigenvalues of $D$ that \cite[lines 1291--1294]{Hastings2006Solving} uses;
    the polar decomposition removes the need for $D$ to be Hermitian.
\end{proof}

\begin{theorem}[The trace-norm form of the cluster-expansion bound]
    \label{thm:hastings_thermal_trace_norm}
    \lean{TNLean.Lattice.ClusterExpansion.re_trace_sqrt_sub_le_of_forall_norm_trace_mul_sub_le}
    \leanok
    \uses{thm:hastings_thermal_trace_norm_duality}
    Let $\rho$ and $\sigma$ be square complex matrices with
    $\operatorname{Re}\operatorname{Tr}\rho\ge0$, and let $\epsilon\ge0$. If
    \begin{align}
        |\operatorname{Tr}(O\rho)-\operatorname{Tr}(O\sigma)|
        \le\epsilon\,\|O\|\operatorname{Re}\operatorname{Tr}\rho
    \end{align}
    for every matrix $O$, then
    $|\rho-\sigma|_1\le\epsilon\operatorname{Re}\operatorname{Tr}\rho$. For the
    thermal state $\rho=\rho(\beta)$, which is positive with
    $\operatorname{Tr}\rho(\beta)=Z=|\rho(\beta)|_1$, and
    $\sigma=\rho(\beta,l_{\mathrm{proj}})$, the hypothesis is the bound of
    eq.~(perc2) with
    $\epsilon=\exp[N\exp(-l_{\mathrm{proj}}/\xi_{\mathrm{perc}})]-1$, and the
    conclusion is the bound of \cite[lines 1281--1299]{Hastings2006Solving}.
\end{theorem}
\begin{proof}\leanok
    For $\|O\|\le1$ the hypothesis gives
    $|\operatorname{Tr}(O(\rho-\sigma))|\le\epsilon\operatorname{Re}\operatorname{Tr}\rho$,
    since $\epsilon\ge0$. Apply
    Theorem~\ref{thm:hastings_thermal_trace_norm_duality} to $D=\rho-\sigma$.
\end{proof}

\begin{theorem}[The matrix product form at every temperature]
    \label{thm:hastings_thermal_all_temperatures}
    \lean{TNLean.Lattice.ClusterExpansion.HasProductForm.pow_of_even,
        TNLean.Lattice.ClusterExpansion.posSemidef_pow_of_even,
        TNLean.Lattice.ClusterExpansion.posSemidef_pow_two_mul,
        TNLean.Lattice.ClusterExpansion.exists_even_div_le}
    \leanok
    \uses{def:hastings_product_form}
    Let $\rho$ be a Hermitian operator in matrix product form of range $l$, and
    let $m$ be even. Then $\rho^m$ is positive semidefinite and in matrix product
    form of range $l$. Moreover, for $\beta\ge0$ and $\beta_0>0$ there is an even
    $m\ge2$ with $\beta/m\le\beta_0$. Applied to the truncated high-temperature
    form $\rho(\beta/m,l_{\mathrm{proj}})$, which is Hermitian, this is the
    positivity and the matrix product form of
    $\rho(\beta,l_{\mathrm{proj}})=\rho(\beta/m,l_{\mathrm{proj}})^m$ in
    eq.~(mplt) of \cite[lines 1303--1313]{Hastings2006Solving}, with positive
    semidefiniteness in place of the source's positive definiteness, which
    would need $\rho(\beta/m,l_{\mathrm{proj}})$ to be
    invertible~\cite{gap:hastings05_mplt_positivity}. The count of
    index values in lines 1313--1320 is not established; it is the count
    recorded as open for
    Theorem~\ref{thm:hastings_trotter_product_form_radius}.
\end{theorem}
\begin{proof}\leanok
    \uses{thm:hastings_product_form_mul}
    For $m=2k$, $\rho^m=(\rho^k)^\dagger\rho^k$ since $\rho^k$ is Hermitian, and
    this is positive semidefinite. Products of operators in matrix product form
    of range $l$ are again in that form, with the indices at each site grouped
    into one (Theorem~\ref{thm:hastings_product_form_mul}). For the number of
    factors take $m=2(\lceil\beta/\beta_0\rceil+1)$.
\end{proof}
